% =====================================================================
%  Thème 1 — Composantes et norme — Niveau DIFFICILE — Corrigé
% =====================================================================
\documentclass[11pt,a4paper]{article}
\usepackage[utf8]{inputenc}
\usepackage[T1]{fontenc}
\usepackage{lmodern}
\usepackage[french]{babel}
\usepackage{amsmath,amssymb,mathtools}
\usepackage{xcolor}
\usepackage{enumitem}
\usepackage{fancyhdr}
\usepackage[a4paper,margin=2.1cm,headheight=26pt,headsep=10pt]{geometry}

\definecolor{primary}{RGB}{222,70,80}
\definecolor{primarydark}{RGB}{150,30,42}
\definecolor{excol}{RGB}{0,135,90}

\pagestyle{fancy}\fancyhf{}
\fancyhead[L]{\small\textcolor{primarydark}{\textbf{Fiche 2} — Calcul vectoriel}}
\fancyhead[R]{\small\textcolor{primarydark}{\textsc{Thème 1 — Difficile — Corrigé}}}
\fancyfoot[C]{\small\thepage}
\renewcommand{\headrulewidth}{0.8pt}

\newcommand{\vc}[1]{\overrightarrow{#1}}
\providecommand{\norm}[1]{\left\lVert #1\right\rVert}
\newcommand{\cor}[1]{\medskip\noindent{\color{primarydark}\bfseries\large Corrigé #1.}\par\smallskip}
\newcommand{\rep}{\textcolor{excol}{$\blacktriangleright$}\ }

\begin{document}
\begin{center}
{\LARGE\bfseries\color{primarydark}Composantes et norme — Corrigé Difficile}\\[2pt]
{\color{primary}\rule{0.5\linewidth}{1.1pt}}
\end{center}
\vspace{2mm}

\cor{1}
\begin{enumerate}[label=\textbf{\alph*)},leftmargin=2em,itemsep=3pt]
  \item $\vc{OP_1}\,(3;4)$, longueur $\sqrt{3^2+4^2}=\sqrt{25}=5$.\\
        $\vc{P_1P_2}\,(3-3\,;\,9-4)=(0;5)$, longueur $\sqrt{0+25}=5$.\\
        $\vc{P_2P_3}\,(-1-3\,;\,12-9)=(-4;3)$, longueur $\sqrt{16+9}=\sqrt{25}=5$.\\
        Les trois trajets mesurent bien $5$~km chacun.
  \item Déplacement net : $\vc{OP_3}\,(-1-0\,;\,12-0)=(-1;12)$, donc
        \[\norm{\vc{OP_3}}=\sqrt{(-1)^2+12^2}=\sqrt{1+144}=\sqrt{145}\approx 12{,}0~\text{km}.\]
  \item Distance \emph{parcourue} $=5+5+5=15$~km. Distance \emph{à vol d'oiseau}
        $=\sqrt{145}\approx 12{,}0$~km. Le drone parcourt plus de chemin que la distance directe
        (les trajets ne sont pas alignés) : c'est logique, un déplacement net est toujours
        $\leq$ à la longueur du chemin suivi.
\end{enumerate}
\rep Le déplacement net se lit directement sur les positions initiale et finale (\og{}extrémité moins
origine\fg{}), sans repasser par les étapes.

\cor{2}
\begin{enumerate}[label=\textbf{\alph*)},leftmargin=2em,itemsep=3pt]
  \item $\vc{MA}\,(1-x\,;2)$ donc $MA^2=(1-x)^2+2^2=(x-1)^2+4$.\\
        $\vc{MB}\,(6-x\,;3)$ donc $MB^2=(6-x)^2+3^2=(x-6)^2+9$.
  \item $MA=MB \iff MA^2=MB^2$ :
        \[(x-1)^2+4=(x-6)^2+9.\]
        On développe : $x^2-2x+1+4=x^2-12x+36+9$, soit $-2x+5=-12x+45$, d'où
        $10x=40$ et $\boxed{x=4}$.
  \item $M\,(4;0)$. Distance commune : $MA=\sqrt{(4-1)^2+2^2}=\sqrt{9+4}=\sqrt{13}$
        (on vérifie $MB=\sqrt{(4-6)^2+3^2}=\sqrt{4+9}=\sqrt{13}$).
\end{enumerate}
\rep Astuce : élever au carré fait disparaître les racines \emph{et} les termes $x^2$ se simplifient,
laissant une équation du premier degré.

\cor{3}
\begin{enumerate}[label=\textbf{\alph*)},leftmargin=2em,itemsep=3pt]
  \item $\norm{\vec u}^2=(\sqrt3-1)^2+(\sqrt3+1)^2=(4-2\sqrt3)+(4+2\sqrt3)=8$, donc
        $\norm{\vec u}=\sqrt8=2\sqrt2$.
  \item $\norm{\vec v}^2=(1+\sqrt2)^2+(1-\sqrt2)^2=(3+2\sqrt2)+(3-2\sqrt2)=6$, donc
        $\norm{\vec v}=\sqrt6$.
\end{enumerate}
\rep Les doubles produits $\pm2\sqrt3$ (ou $\pm2\sqrt2$) se \textbf{compensent} quand on additionne
les carrés de deux quantités conjuguées : le résultat est un entier.

\cor{4}
\begin{enumerate}[label=\textbf{\alph*)},leftmargin=2em,itemsep=3pt]
  \item $\vc{AB}\,(t-1\,;\,(t+3)-2)=(t-1\,;\,t+1)$.
  \item $\norm{\vc{AB}}^2=(t-1)^2+(t+1)^2=(t^2-2t+1)+(t^2+2t+1)=2t^2+2$.
  \item $2t^2+2$ est minimal quand $t^2$ est minimal, c'est-à-dire pour $t=0$.
        Alors $\norm{\vc{AB}}^2=2$ et $AB=\sqrt2$.
  \item $AB=\sqrt{10}\iff 2t^2+2=10\iff t^2=4\iff t=2 \text{ ou } t=-2$.
\end{enumerate}
\rep Travailler avec $\norm{\vc{AB}}^2$ évite les racines : c'est une fonction du second degré en $t$,
minimale au sommet ($t=0$ ici).

\end{document}
